\title{LongRoPE: Extending LLM Context Window Beyond 2 Million Tokens}

\begin{abstract}
A substantial context window represents a valuable capability in large language models (LLMs). Nonetheless, existing methods for extending context length are constrained—typically not exceeding 128k tokens—because of the substantial expense of fine-tuning, the limited availability of sufficiently lengthy documents, and the disruptive effects caused by introducing novel token positions.
\end{abstract}

This work presents LongRoPE, a novel approach that, for the first time, expands the context window of existing pre-trained LLMs to an extraordinary \textbf{2048k} tokens, requiring as few as 1k fine-tuning steps using training sequences no longer than 256k tokens, all while preserving performance at the original, shorter context window. This advancement stems from three primary contributions: (i) we discover and leverage two types of non-uniformities within positional interpolation via an efficient search process, yielding superior fine-tuning initialization and facilitating an 8$\times$ context extension even without fine-tuning; (ii) we propose a staged extension methodology in which a 256k-token LLM is first fine-tuned, after which a second round of positional interpolation on this extended model enables achieving a context window of 2048k tokens; (iii) we recalibrate LongRoPE on an 8k-token length setting to restore optimal performance at shorter context windows. Comprehensive evaluations using LLaMA2 and Mistral on diverse tasks substantiate the method’s efficacy. LongRoPE-extended models preserve their original architecture, involving only minor adjustments to positional embeddings, and are compatible with most existing optimizations. Code will be released at \texttt{https://github.com/microsoft/LongRoPE}

\section{Introduction}

Although Large Language Models (LLMs) have demonstrated impressive capabilities across a wide range of tasks~\cite{gpt4,llama2}, their performance is constrained by a restricted context window—for example, LLaMA2 is limited to 4096 tokens~\cite{llama2}. When input exceeds this window, LLMs' effectiveness drops, as they are confronted with positional encodings beyond what they encountered during training. This restriction leads to practical difficulties in key settings such as in-context learning with many exemplars~\cite{Huang2023FewerIM} and in the operation of LLM agents~\cite{park2023generative,madaan2023selfrefine}.

Recent studies have demonstrated that it is possible to expand the context window of pre-trained LLMs to approximately 128k through fine-tuning with longer input sequences~\cite{longlora,pi,yarn,zhang2024soaring,liu2023scaling}. Nevertheless, three principal challenges hinder the extension of the context window beyond this limit. First, introducing untrained position indices results in numerous aberrant values, thereby causing distributional shifts and hampering the convergence of fine-tuning~\cite{pi}. This becomes especially problematic when the context length is increased from 4k to over 1000k, which involves incorporating more than 90\% novel positional embeddings. Second, effective fine-tuning typically depends on the availability of texts matching the extended lengths. Yet, there is a paucity of long documents—particularly those surpassing 1000k tokens—in existing datasets. Furthermore, training with these extra-long sequences is highly demanding computationally, as it consumes an excessive amount of GPU resources and time. Third, as the context window is greatly expanded, the attention mechanism’s focus is diluted across a vast array of token positions, leading to a decline in performance on shorter sequences that the model previously handled well~\cite{pi}.

A commonly used strategy to address the initial challenge involves the interpolation of RoPE positional embeddings~\cite{rope,pi}, which adjusts the positional indices of new tokens so they fit within the original range used during pretraining (see Fig.\ref{fig:overview}). The Position Interpolation (PI) technique~\cite{pi} achieves this by performing a linear interpolation on the rotary angles of RoPE according to the scaling ratio. Alternatively, NTK~\cite{ntk,dynamicntk} introduces non-uniform interpolation and extrapolation schemes applied across different RoPE dimensions. YaRN~\cite{yarn} further segments the RoPE dimensions into three frequency-defined categories, assigning extrapolation, NTK-based interpolation, and linear interpolation methods to each respective group. Nevertheless, the information entropy captured by positional embeddings within the Transformer model reveals \emph{intricate and non-uniform distribution patterns}. Current methods fail to adequately exploit this subtle non-uniformity, which results in the loss of informational content and imposes constraints on the expansion of the context window.

Section~\ref{sec:analysis} provides empirical evidence for two major insights: \textbf{(1)} Successful positional interpolation necessitates accounting for two types of non-uniformities: differences across RoPE dimensions and across token positions. Minimal interpolation should be applied to lower RoPE dimensions and to earlier token positions, but the ideal interpolation strategy is contingent on the target extension length.
  \textbf{(2)} Incorporating these non-uniform factors into the positional interpolation approach enables preservation of crucial information inherent in the original RoPE, particularly in key dimensions and for certain token positions. This strategy reduces information loss resulting from interpolation and thus yields superior initialization when subsequently fine-tuning. Additionally, this method supports extension by up to 8$\times$ even without fine-tuning.

Building upon these insights, we present \textbf{LongRoPE}, a robust approach designed to expand the LLM context window to lengths exceeding 2 \emph{million} tokens. The design of LongRoPE rests on three principal innovations. Firstly, {LongRoPE} leverages the multidimensional and non-uniform nature of positional interpolation, determining suitable rescale factors for RoPE's rotation angles across each RoPE dimension based on token positions. Given that the complexity of searching for these rescale factors rises exponentially as the targeted extension ratio increases, {LongRoPE} employs an evolutionary search algorithm augmented by two specific optimization strategies to enhance the efficiency of the search process. An illustration of the rescaled RoPE identified through this procedure is depicted in Fig.~\ref{fig:overview}.

Subsequently, {LongRoPE} utilizes a resource-effective, incremental extension approach to reach a 2048k context window, avoiding the necessity for direct fine-tuning on ultra-long text sequences, which are both rare and difficult to procure. This approach starts by identifying a suitable 256k sequence length for the pre-trained LLM and proceeds with fine-tuning at this scale. Following this, taking advantage of the non-uniform positional interpolation’s capacity for an 8$\times$ extension in scenarios without fine-tuning, a second stage involves searching for updated RoPE rescale factors on the already fine-tuned, extended LLM. Through these steps, a 2048k context window is ultimately realized for both LLaMA2 and Mistral~\cite{mistral}.

In order to address potential drops in performance on the original, shorter context window, {LongRoPE} further tunes the RoPE rescaling parameters in the expanded LLM. Analogous to the process of scaling context from 256k up to 2048k, the model is also calibrated downward to 4k and 8k context windows using our optimization algorithm, aiming to reduce excessive positional interpolation. At inference time, when the input sequence length is under 8k, RoPE is modified with the previously determined rescale parameters.

Comprehensive experimentation with multiple LLM architectures and a range of long-context tasks validates the efficacy of our approach. Our findings indicate that LongRoPE consistently preserves low perplexity across evaluation lengths from 4k up to 2048k, surpasses 90\% passkey retrieval accuracy, and attains accuracy on par with conventional benchmarks situated within a 4096-token context window. The LongRoPE technique is adaptable to any LLM utilizing RoPE embeddings. We intend to make both our codebase and the LongRoPE-2048k models publicly available.

\section{Non-uniformity in Positional Interpolation}
\label{sec:analysis}

\subsection{Preliminary}

Transformer architectures necessitate the incorporation of positional data—typically achieved via position embeddings—to encode token order within sequences. In this study, we center our attention on the RoPE~\cite{rope} position embedding scheme, which has become prevalent in state-of-the-art LLMs. The RoPE embedding associated with a token located at position $n$ can be concisely expressed as:

\begin{equation}
		\fontsize{7.8}{7.8} \selectfont
	\label{eq:rope}
	\left[ cos(n\theta_0), sin(n\theta_0), cos(n\theta_1), \cdot\cdot\cdot, cos(n\theta_{d/2-1}), sin(n\theta_{d/2-1}) \right]   
\end{equation}

Here, $d$ stands for the embedding dimension, while $n\theta_i$ denotes the rotary angle associated with the token located at position $n$. The term $\theta_i=\theta^{-2i/d}$ specifies the corresponding rotation frequencies. For RoPE, $\theta$ typically takes a default value of 10000.

\textbf{\emph{Context window extension ratio $s$} and positional interpolation}. We define $s$ as  the ratio of extended context length $L'$ to the original length $L$: $s=\frac{L'}{L}$. 

To increase the context window from $L$ to $L'$, contemporary positional interpolation techniques propose reducing the rotation frequencies $\theta_i$ proportionally to the extension ratio $s$. Given that $\beta = \theta^{2/d}$ and $\lambda$ represents the corresponding scaling factor determined by $s$, we represent these positional interpolation strategies in a unified manner as:

\begin{equation}	
	\fontsize{7.5}{7.5} \selectfont
	\label{eq:unified_rope}
	\begin{aligned}
		\left[cos\left(\frac{n}{\lambda(\beta)^0}\right), sin \left(\frac{n}{\lambda(\beta)^0}\right), cos\left(\frac{n}{\lambda(\beta)^1}\right), \cdot\cdot\cdot,  sin \left(\frac{n}{\lambda(\beta)^{d/2-1}}\right) \right]   
	\end{aligned}
\end{equation}

\textbf{Linear positional interpolation (PI)}. PI~\cite{pi} introduces a method where position indices are linearly interpolated so that they fit within the pre-trained positional range. When the sequence is extended by a factor of $s$, each RoPE dimension's rotational angle is proportionally compressed by a factor of $\lambda=s$. While this approach enables handling longer contexts, it causes the positional encodings to be overly compressed, making it more difficult for the model to tell apart nearby tokens. As a result, PI generally exhibits poorer performance at larger extension ratios.

\textbf{NTK-based interpolation and extrapolation}. The studies in~\cite{ntk,dynamicntk} approach RoPE through the lens of information representation, leveraging the framework of Neural Tangent Kernel (NTK) theory~\cite{jacot2018neural,tancik2020fourier}. To address the problem of crowded positions inherent in PI, their method redistributes the interpolation demand more evenly across the various RoPE dimensions. This strategy entails scaling the lower (higher-frequency) dimensions to a lesser degree, while scaling higher (lower-frequency) dimensions more significantly. Such an approach facilitates both positional interpolation and extrapolation using $\lambda=s^{i}$. The enhanced dynamic NTK method~\cite{dynamicntk} further adapts the scaling ratio at each position in accordance with the current length of the sequence. In contrast to PI, which relies on additional fine-tuning, approaches informed by NTK can broaden context window size even without fine-tuning, albeit with a typical maximum extension ratio of 4$\times$.

\textbf{YaRN}~\cite{yarn} proposes a notable advancement in the effectiveness of positional interpolation. The method organizes RoPE dimensions according to their frequencies into three distinct categories, assigning a specific interpolation approach to each group. Extrapolation ($\lambda$=1) is applied to high-frequency dimensions, linear interpolation (PI) is designated for those with low frequencies, and dimensions with intermediate frequencies utilize the NTK technique. The central feature of YaRN is this partitioning of RoPE dimensions, which currently relies on empirical grouping guided by human judgment. As a consequence, its applicability to new LLMs might lead to less-than-optimal outcomes.

\subsection{Study on Non-uniform Positional Interpolation}

Building on insights from NTK and YaRN, we observe that their improvements derive from incorporating non-linear elements, particularly by handling various frequency components separately within RoPE dimensions for tailored interpolation and extrapolation tasks. Nonetheless, the dominant approaches to introducing non-linearity depend largely on heuristics crafted by researchers. This leads us to two central inquiries: (1) Does the prevailing approach to positional interpolation represent the best possible method? (2) Might there exist alternative forms of non-linearity that have yet to be investigated?

In order to address these questions, we apply evolution search (refer to Sec.~\ref{sec:method}) to identify improved non-uniform positional interpolations tailored for LLaMA2-7B. This search process is driven by perplexity scores, utilizing 5 randomly chosen examples from the PG19~\cite{pg19} validation dataset. Our empirical investigation leads us to highlight several important findings.

\textbf{Finding 1}: \textit{RoPE dimensions exhibit substantial non-uniformities, which are not effectively handled by current positional interpolation methods.}

We optimize $\lambda$ individually for each RoPE dimension as described in Eq.~\ref{eq:unified_rope}. In Table~\ref{tbl:exp1}, we present a perplexity comparison of LLaMA2-7B across different interpolation techniques on the PG19 and Proof-pile~\cite{proof-pile} test datasets, all evaluated in a zero-shot (non-fine-tuned) setting. The outcomes from our optimized approach demonstrate marked perplexity reductions, indicating that prevalent strategies such as linear (PI) and non-uniform (Dynamic-NTK and YaRN) interpolation do not achieve optimal performance. Interestingly, YaRN performs worse than both PI and NTK on PG19, owing to its failure to effectively extend to the intended context window size in non-fine-tuned LLMs. For instance, on an 8k context, YaRN exhibits a steep increase in perplexity beyond 7k tokens.

Our investigation reveals that the rescaling factors $\lambda$ in Eq.~\ref{eq:unified_rope} exhibit variability, in contrast to the constant scaling factor $s$ used in PI, NTK's formulaic approach, and YaRN's group-based method. The application of these variable, non-uniform scaling factors leads to a marked enhancement in LLaMA2's language modeling capabilities—specifically in terms of perplexity—when evaluated on 8k and 16k context lengths, all without requiring any additional fine-tuning. This improvement is attributed to the fact that the resultant positional embedding more faithfully retains characteristics of the original RoPE, particularly in essential dimensions, which alleviates the model’s challenge in discerning nearby token positions.

\textbf{Finding 2}: \textit{RoPE for the initial tokens in the input sequence should be extrapolated with less interpolation.} 

We propose that for the first $\hat{n}$ tokens in input sequences, reduced interpolation should be applied in their RoPE encoding. This stems from their substantial attention scores—an observation consistent with findings from Streaming LLM~\cite{streamingllm} and LM-Infinite~\cite{han2023lminfinite}—which suggest their prominent role within attention mechanisms. To empirically test this idea, we conduct experiments expanding the context window to 8k and 16k using both PI and NTK, but omit interpolation for the initial $\hat n$ tokens ($\hat n = 0, 2, \ldots, 256$). For $\hat n = 0$, the scheme defaults to standard PI or NTK methods. Table~\ref{tbl:tokenposition} presents two main findings: \textbf{(1)} excluding position interpolation for the initial tokens consistently enhances performance, and \textbf{(2)} the ideal value for $\hat n$ is contingent on the target window length.

\textbf{Finding 3}: \textit{Non-uniform positional interpolation effectively extends LLM context window in both fine-tuning and non-fine-tuning settings.} 

Our investigation demonstrates that implementing our optimized non-uniform position interpolation method leads to notable enhancements in extension performance at 8k and 16k context lengths without requiring fine-tuning. However, extending to longer context windows still necessitates additional fine-tuning. To address this, we apply fine-tuning to LLaMA2-7B using our optimized RoPE approach for a 64k context length (refer to the Appendix for fine-tuning details). As reported in Table~\ref{tbl:finetune}, our approach consistently surpasses both PI and YaRN, regardless of whether LLaMA2-7B has been fine-tuned. This performance gain is attributed to our strategic exploitation of non-uniform positional interpolation, which minimizes information degradation and furnishes a superior initialization for the fine-tuning stage.

\textbf{Summary}. Our research identifies two distinct non-uniformities: differences across RoPE dimensions and across token positions. Harnessing these non-uniformities in positional interpolation yields substantial gains in extending the context windows of LLMs.

\section{LongRoPE}

\label{sec:method}
Building on these insights, we propose LongRoPE, a method that initially incorporates an optimized search algorithm to take full advantage of both types of non-uniformity. Leveraging this approach, LongRoPE extends the context window of LLMs to surpass 2 million tokens.

\subsection{Problem Formulation}

These two forms of non-uniformity result in an extensive solution space and add layers of complexity to the optimization process. To mitigate these challenges, we reformulate the problem of multidimensional non-uniform position interpolation optimization into a search framework.

Consider a large language model (LLM) designed for a context window of size $L'$ and a set of long-form documents $\mathbf{X}$, where each document $\mathbf{x}\in\mathbf{X}$ contains more tokens than $L'$. The standard rotary angle for the $i^{th}$ dimension in RoPE embedding at token position $n$ is represented as $\frac{n}{\beta^i}$. This gives rise to the following optimization problem:

\begin{equation}
\begin{gathered}
		\label{eq:problem}

\underset {\mathbf{x}\in\mathbf{X}; \, |\mathbf{x}|\ge L'}{ \operatorname {arg\,min} } \,  \mathcal{L}\, \left( \text{LLM} (\text{RoPE}, \mathbf{X})\right),\text{where}
\\
\scalebox{0.93}{
$\underset {\begin{subarray}{l} i=0,\ldots,\frac{d}{2}-1;\\ n\in[0,|\mathbf{x}|);\end{subarray}}{\text{RoPE}_(n)}= {\left[\cdots, \cos\left(\mathbb{I}(\hat{\lambda}_{i}, \hat{n}) \cdot \frac{n}{\beta^i}\right),  \sin\left(\mathbb{I}(\hat{\lambda}_{i}, \hat{n}) \cdot \frac{n}{\beta^i}\right), \cdots\right]}$}
\\
\text{where} \;
\mathbb{I}(\hat{\lambda}_{i}, \hat{n}) = \begin{cases}
    1 & \text{if $n<\hat{n}$} \\
    \frac{1}{\lambda_{i}} & \text{if $n \geq \hat{n}$}
\end{cases}
\end{gathered}
\end{equation}

In this context, we define a collection of rescaling factors, $\mathbb{I}(\hat{\lambda}_{i},\hat{n})$, to account for both types of non-uniformities. Here, $\hat{\lambda_i}$ characterizes the non-uniformity present in RoPE dimensions, while $\hat{n}$ reflects positional non-uniformity among tokens. The function $\mathbb{I}(\hat{\lambda}_{i},\hat{n})$ is utilized to adjust the rotational angle corresponding to the $i^{th}$ RoPE dimension, such that $\hat\lambda_{i}$ serves as the adjustment parameter and $\hat{n}$ defines the token position cutoff. For the first $\hat{n}-1$ token indices, no scaling is applied and the standard RoPE rotation angle $\frac{n}{\beta^i}$ remains unchanged. However, when $n\ge\hat{n}$, the rescaling factor is incorporated.

With a specified target context window length of $L'$, our goal is to determine the set of optimal rescaling factors ($\mathbb{I}(\hat{\lambda}_{0},\hat{n})$, $\mathbb{I}(\hat{\lambda}_{1},\hat{n})$, ..., $\mathbb{I}(\hat{\lambda}_{i},\hat{n})$, ...) across all RoPE dimensions, ranging from the $1^{st}$ to the $d^{th}$. Consequently, employing the rescaled $\text{RoPE}$ within the target $\text{LLM}$ enables the model to minimize the next token prediction loss $\mathcal{L}$—that is, the perplexity—when processing input sequences $\mathbf{X}$ consisting of $L'$ tokens.

\subsection{Searching the Non-uniform Position Interpolation}

To address the challenge posed in Eq.~\ref{eq:problem}, we propose an efficient and straightforward approach that seeks the best RoPE rescaling factors, thereby maximizing the utilization of multidimensional variations present in position embeddings.

\textbf{Search space}. We construct an extensive search space that encompasses both forms of non-uniformity. Table~\ref{tbl:searchspace} provides an overview of this space. In particular, our strategy enables the assignment of distinct rescale factors to each dimension in the RoPE embedding. To simplify the design of the search space, we opt to directly optimize over $\lambda_i$ and $\hat{n}$, rather than the indicator $\mathbb{I}(\hat{\lambda}_{i},\hat{n})$, noting that $\hat{\lambda}_{i}=1/\lambda_i$. As indicated in Table~\ref{tbl:searchspace}, $\lambda_i$ can take values beginning at 1.0 (corresponding to straightforward extrapolation) and extending up to $s\times1.25$ (representing more aggressive interpolation than PI), with increments of 0.01, where $s$ denotes the intended expansion ratio of the context window.

The parameter $\hat{n}$ specifies how many of the initial token positions maintain their original configuration by bypassing position interpolation, instead utilizing the standard RoPE embedding. In practice, we select $\hat{n}$ from the set \{0, 1, 2, 4, 8, 12, 16, 20, 24, 28, 32, 64, 128, 256\}. If $\hat{n}=0$, then every token position is assigned the corresponding rescale factors found through the search process.

\textbf{Evolution-based search}. The search space detailed in Table~\ref{tbl:searchspace} encompasses a vast array of positional interpolation configurations, which presents considerable difficulties in navigating the space efficiently. As an example, using an extension factor of $s=4\times$ results in $400^{128/2}\times14$=4$\times10^{167}$ possible options. When the extension ratio increases, the search space size grows exponentially. To efficiently handle this complexity, we adopt evolution search~\cite{guo2020single} and propose two optimization strategies that substantially enhance the efficiency of the search process. The overall procedure is summarized in Algorithm~\ref{alg:search}.

\textit{Optimized initial population generation}. Rather than relying solely on random initialization for the $P$ rescale factors, the three RoPE rescale factors associated with PI, NTK, and YaRN are explicitly included as members of the initial population. The rest of the $P$-3 individuals are produced by applying random mutations to these three rescale factors, each with a mutation probability of $p$.

\textit{Monotonically non-decreasing constraint}. Following the creation of the initial population, we assess each individual by calculating the LLM perplexity. This involves applying the appropriate RoPE rescale factors within the target LLM and then measuring the perplexity on the input $\mathbf{X}$. The $k$ individuals with the lowest perplexity are then selected as parents for the evolutionary process. Nonetheless, because the search space is extremely large, unrestrained mutation and crossover operations may frequently encounter low-quality solutions, which results in excessive and redundant perplexity evaluations. This inefficiency is exacerbated for large $L'$, as each perplexity inference requires considerable computational resources.

To mitigate this issue, we enforce a monotonicity constraint such that the sampled RoPE rescaling factors are non-decreasing, i.e., $\lambda_i \le \lambda_{i+1}$. Only those RoPE configurations fulfilling this criterion are utilized in the LLM’s perplexity assessment, thereby curtailing the computational overhead of the search process. In particular, we stipulate that $\lambda_i$ grows monotonically as the RoPE dimension index $i$ (where $i=0,\ldots,63$) increases. This requirement for dimensional monotonicity draws from NTK theory~\cite{jacot2018neural,tancik2020fourier,ntk}, which posits that dimensions with higher frequency components (lower $i$) are less dependent on interpolation—necessitating smaller $\lambda_i$ values—whereas dimensions associated with lower frequencies (higher $i$) can sustain larger degrees of interpolation, thus supporting larger $\lambda_i$ values.

\begin{algorithm}[t]
	\small
	\caption{The search algorithm}
	\textbf{Input:} target LLM, input samples $\mathbf{X}$, population size $P$, mutation size $N_1$, crossover size $N_2$, max iterations $\mathcal{T}$, mutate probability $p$\\

	\begin{algorithmic}[1]
		\label{alg:search}
		\STATE Top-k$\leftarrow \phi$;
		\STATE $\text{P}_0 \gets \textit{Initialize\_population\_with\_optimization}(P, \mathbf{X}, p)$;
		\FOR{$i \gets 1$ to $\mathcal{T}$}
		\STATE \textit{Compute\_perplexity}(LLM, $\text{P}_{i-1}$, $\mathbf{X}$);
		\STATE Top-k $\gets$ \textit{Update\_Topk}(Top-k, $\text{P}_{i-1}$);
		\STATE $\text{P}_{mutation} \gets \textit{Mutation\_with\_mono\_constraint}$(Top-k, $N_1$, $p$);
		\STATE $\text{P}_{crossover} \gets \textit{Crossover\_with\_mono\_constraint}$(Top-k, $N_2$);
		\STATE $\text{P}_i \gets \text{P}_{mutation} \cup \text{P}_{crossover} \cup$ Top-k;
		\ENDFOR
		\STATE Output the element in Top-k achieving minimum perplexity;
	\end{algorithmic}
\end{algorithm}

\textbf{8$\times$ extension without fine-tuning}. Utilizing evolutionary search, our approach efficiently discovers non-uniform RoPE rescale factors, safeguarding essential dimensions and positions in order to reduce the information degradation caused by interpolation. As illustrated in Fig.\ref{fig:non-ft}, this method successfully expands LLaMA2’s context window from 4k up to 32k without requiring any fine-tuning. By comparison, other techniques like PI and the non-uniform versions of NTK and YaRN experience a significant increase in perplexity following a 2$\times$ extension.

\subsection{Extending LLM Context Window to 2048K}
\label{sec:extendto2048k}

\textbf{Progressive extension to 2048k}. In this section, we present our approach for expanding the context window of pre-trained LLMs from the standard 4k up to 2048k tokens. Our non-uniform positional interpolation strategy enables an 8$\times$ increase in context length without requiring additional fine-tuning. When aiming for more substantial expansions (such as 512$\times$), however, fine-tuning becomes essential. A possible technique is to determine optimal RoPE rescaling factors suitable for the desired 2048k window, followed by fine-tuning the model. Nevertheless, this procedure entails significant computational costs, making it resource-intensive. Furthermore, in our experience, effectively fine-tuning LLMs at such high extension ratios presents considerable difficulties (refer to Appendix for details).

Fortunately, LongRoPE demonstrates efficacy when applied to both the base and the fine-tuned extended LLM. Thus, we propose an efficient and incremental approach that enables reaching the desired 2048k context length using only 1k fine-tuning steps, each performed within a training window of 256k tokens.

$\Diamond$ \textit{Exploring 256k context window expansion of pre-trained LLM using LongRoPE search}. Using LLaMA2 for demonstration, we perform a search process aimed at achieving context window lengths of 128k and 256k tokens. At this phase, the expansion ratios are set to 32$\times$ and 64$\times$, accordingly.

$\Diamond$ \textit{Fine-tuning to 256k}. In the subsequent phase, we adapt the pre-trained LLM to extend its context window to 256k. Initially, we fine-tune LLaMA2 for 400 steps utilizing the RoPE rescaled factors set for 128k. After this stage, we update the model by swapping in the RoPE rescaled factors for 256k at the resultant checkpoint, followed by a further 600 steps of fine-tuning. This staged approach has been found to be more efficient compared to directly fine-tuning the model to 256k.

$\Diamond$ \textit{Adapting fine-tuned extended LLM to 2048k using LongRoPE search}. As a concluding step, we conduct an additional search process utilizing the fine-tuned LLM at a 256k sequence length. This method enables an expansion to a remarkably large context window of 2048k tokens without necessitating any further fine-tuning procedures. The achieved overall extension ratio is $512\times$.

\textbf{Recovery within limited context windows}. Upon increasing the context window length to a substantial 2048k tokens, a decline in model effectiveness is observed when evaluating the initial context window. This phenomenon is well-documented with positional interpolation~\cite{pi}, as enlarging the context causes high-dimensional position embeddings in the base context range to be confined to a compressed subspace, impairing the model’s capabilities. In the scenario of a 512$\times$ extension, the position representations inside the baseline 4k window experience significant compression.

To address this issue, we conduct an additional evolution-based search on the extended LLM to refine the RoPE rescale factors specifically for shorter context lengths such as 4k and 8k. For these shorter sequences, we lower the upper limit of the searched $\lambda$ values since minimal positional interpolation is necessary. At inference time, the LLM automatically tunes the appropriate RoPE rescale factors as needed.

\section{LongRoPE}

\label{sec:method}
Building upon these insights, we propose LongRoPE. This approach begins by developing an optimized search algorithm designed to capitalize on the two forms of non-uniformity. Utilizing this algorithm, LongRoPE is able to expand the LLM’s context window to exceed 2 million tokens.

\subsection{Problem Formulation}

These two forms of non-uniformity expand the solution space considerably, resulting in increased complexity for optimization processes. To manage this challenge, we reformulate the multidimensional non-uniform position interpolation optimization issue as a search problem.

Consider a LLM designed for a context window of size $L'$, processing long input documents $\mathbf{X}$, with each segment $\mathbf{x}\in\mathbf{X}$ containing more tokens than $L'$. We represent the standard rotary angle for the $i^{th}$ dimension at token position $n$ within the RoPE embedding as $\frac{n}{\beta^i}$. The corresponding optimization problem can be formally specified as:

\begin{equation}
\begin{gathered}
		\label{eq:problem}

\underset{\mathbf{x}\in\mathbf{X};\, |\mathbf{x}|\ge L'}{\operatorname{arg\,min}}\,\mathcal{L}\left(\text{LLM}\left(\text{RoPE}, \mathbf{X}\right)\right),\text{ where}
\\
\scalebox{0.93}{
$\underset{\begin{subarray}{l}i=0,\ldots,\frac{d}{2}-1;\\ n\in [0,|\mathbf{x}|);\end{subarray}}{\text{RoPE}_{(n)}} = \left[\ldots, \cos\left(\mathbb{I}(\hat{\lambda}_i, \hat{n}) \cdot \frac{n}{\beta^i}\right), \sin\left(\mathbb{I}(\hat{\lambda}_i, \hat{n}) \cdot \frac{n}{\beta^i}\right), \ldots \right]$}
\\
\text{with}\quad \mathbb{I}(\hat{\lambda}_{i},\hat{n}) = \begin{cases}
1 &\text{if } n<\hat{n}\\
\frac{1}{\lambda_i} &\text{if } n\geq \hat{n}
\end{cases}
\end{gathered}
\end{equation}

In this framework, we define a collection of scaling coefficients, $\mathbb{I}(\hat{\lambda}_{i},\hat{n})$, to account for two distinct types of non-uniformities. Here, $\hat{\lambda_i}$ represents irregularities across RoPE dimensions, and $\hat{n}$ indicates variations associated with token positions. More concretely, $\mathbb{I}(\hat{\lambda}_{i},\hat{n})$ serves to adjust the rotation angle assigned to the $i^{th}$ dimension of RoPE, with $\hat\lambda_{i}$ functioning as the scaling coefficient and $\hat{n}$ denoting the position threshold for tokens. For the first $\hat{n}$-1 positions, the standard RoPE rotation angle $\frac{n}{\beta^i}$ is maintained, and the scaling coefficient $\hat\lambda_{i}$ is not applied. In contrast, for all subsequent positions where $n\ge\hat{n}$, the scaling factor is introduced into the computation.

Assuming a desired context window length of $L'$, our aim is to determine the most suitable rescale factors ($\mathbb{I}(\hat{\lambda}_{0},\hat{n})$, $\mathbb{I}(\hat{\lambda}_{1},\hat{n})$, ...$\mathbb{I}(\hat{\lambda}_{i},\hat{n})$...) for each RoPE dimension from the $1^{st}$ up to the $d^{th}$. By employing these rescaled $\text{RoPE}$ factors in the target $\text{LLM}$, the system seeks to minimize the next token prediction loss, denoted by $\mathcal{L}$ (equivalently, the perplexity), when processing input sequences $\mathbf{X}$ that have a token count of $L'$.

\subsection{Searching the Non-uniform Position Interpolation}

To address the challenge posed in Eq.~\ref{eq:problem}, we propose an efficient and straightforward approach that identifies optimal RoPE rescaling factors, thereby maximizing the benefits from the inherent multidimensional non-uniformities present in position embeddings.

\textbf{Search space}. Our method establishes a broad search space that accounts for the two main types of non-uniformities. The construction of this space is detailed in Table~\ref{tbl:searchspace}. Notably, we permit the optimization of distinct rescale factors for every dimension within the RoPE embedding. For simplicity in the search space formulation, our search is conducted over $\lambda_i$ and $\hat{n}$ rather than directly over $\mathbb{I}(\hat{\lambda}_{i},\hat{n})$, with $\hat{\lambda}_{i}$ defined as $1/\lambda_i$. As presented in Table~\ref{tbl:searchspace}, the permissible values for $\lambda_i$ range from 1.0 (corresponding to direct extrapolation) up to $s\times1.25$ (encompassing more extensive interpolation than what is offered by PI), incremented in steps of 0.01; here, $s$ represents the ratio by which the context window is expanded.

The parameter $\hat{n}$ specifies how many of the earliest token positions preserve the standard RoPE embedding, foregoing position interpolation. In practice, $\hat{n}$ is chosen from the set \{0, 1, 2, 4, 8, 12, 16, 20, 24, 28, 32, 64, 128, 256\}. Setting $\hat{n}=0$ results in applying the optimized rescale factors to every token position.

\textbf{Evolution-based search}. The search space outlined in Table~\ref{tbl:searchspace} contains a vast array of potential positional interpolation configurations, making comprehensive search highly demanding. As an illustration, employing an extension of $s=4\times$ results in $400^{128/2}\times14 = 4\times10^{167}$ possible combinations. This search space increases exponentially with larger extension ratios. To cope with this complexity, we adopt an evolution search strategy~\cite{guo2020single} and implement two optimization methods to substantially enhance the efficiency of the search. The full search workflow is depicted in Algorithm~\ref{alg:search}.

\textit{Optimized initial population generation}. Rather than populating the initial set of $P$ rescale factors purely at random, we explicitly include the RoPE rescale factors for PI, NTK, and YaRN as specific members of the initial population. The other $P-3$ individuals are produced by randomly applying mutations to these three baseline rescale factors, each with probability $p$.

\textit{Monotonically non-decreasing constraint}. Once the initial population has been produced, LLM perplexity is evaluated for all individuals. This involves applying the respective RoPE rescale factors within the target LLM and subsequently calculating the perplexity on input $\mathbf{X}$. The individuals with the top-k performance are selected as parents for the evolutionary phase. Nonetheless, the immense breadth of the search space allows indiscriminate mutation and crossover to sample suboptimal candidates, often resulting in superfluous perplexity evaluations. This inefficiency is particularly pronounced for large values of $L'$, as each perplexity computation requires substantial inference time.

To address this issue, we enforce a monotonic non-decreasing constraint on the rescaled RoPE factors such that $\lambda_i \leq \lambda_{i+1}$. Only those RoPE parameterizations meeting this condition are considered for LLM perplexity assessment, thereby significantly curtailing the computational cost associated with the search. In particular, our approach mandates that $\lambda_i$ increases monotonically across RoPE dimensions, indexed by $i=0,\ldots,63$. The rationale for enforcing monotonicity by dimension is grounded in NTK theory~\cite{jacot2018neural,tancik2020fourier,ntk}, which posits that lower dimensions, associated with higher frequencies, necessitate less interpolation (i.e., smaller $\lambda_i$), whereas higher dimensions at lower frequencies benefit from increased interpolation (i.e., larger $\lambda_i$).

\begin{algorithm}[t]
	\small
	\caption{The search algorithm}
	\textbf{Input:} target LLM, input samples $\mathbf{X}$, population size $P$, mutation size $N_1$, crossover size $N_2$, max iterations $\mathcal{T}$, mutate probability $p$\\

	\begin{algorithmic}[1]
		\label{alg:search}
		\STATE Top-k $\gets \phi$;
		\STATE $\text{P}_0 \gets$ \textit{Initialize\_population\_with\_optimization} ($P$, $\mathbf{X}$, $p$);
		\FOR{$i = 1$ \textbf{to} $\mathcal{T}$}
		\STATE \textit{Compute\_perplexity} (LLM, $\text{P}_{i-1}$, $\mathbf{X}$);
		\STATE Top-k $\gets$ \textit{Update\_Topk} (Top-k, $\text{P}_{i-1}$);
		\STATE $\text{P}_{mutation} \gets$ \textit{Mutation\_with\_mono\_constraint} (Top-k, $N_1$, $p$);
		\STATE $\text{P}_{crossover} \gets$ \textit{Crossover\_with\_mono\_constraint} (Top-k, $N_2$);
		\STATE $\text{P}_i \gets \text{P}_{mutation} \cup \text{P}_{crossover} \cup$ Top-k;
		\ENDFOR
		\STATE Output the member of Top-k with the minimum perplexity;
	\end{algorithmic}
\end{algorithm}

\textbf{8$\times$ extension without fine-tuning}. Through the use of evolutionary search, our approach successfully determines optimal non-uniform RoPE rescale factors that safeguard critical dimensions and positions, thereby reducing information degradation from interpolation. As illustrated in Fig.\ref{fig:non-ft}, this enables the expansion of LLaMA2's context window from 4k to 32k tokens without the need for fine-tuning. By comparison, previous approaches—including PI, as well as non-uniform NTK and YaRN—experience a sharp increase in perplexity beyond a 2$\times$ extension.

\subsection{Extending LLM Context Window to 2048K}
\label{sec:extendto2048k}

\textbf{Progressive extension to 2048k}. In this section, we present our approach for scaling the context window of pre-trained LLMs from the standard 4k tokens up to over 2048k. Our technique, which leverages non-uniform positional interpolation, enables an 8$\times$ increase in context length without the need for fine-tuning. However, achieving more substantial extensions (such as 512$\times$) does necessitate fine-tuning. One possible strategy involves identifying optimal RoPE rescaling factors tailored for the 2048k target and conducting fine-tuning afterward. Nonetheless, this approach is confronted by the extremely high cost of training at such scales. Furthermore, our observations indicate that effectively fine-tuning LLMs at such significant extension ratios remains difficult (refer to Appendix for further details).

Fortunately, LongRoPE demonstrates efficacy with both base and fine-tuned extended LLMs. As a result, we propose a streamlined, incremental approach that attains the desired 2048k context window using merely 1k fine-tuning iterations at a training length capped at 256k.

$\Diamond$ \textit{Searching for extending pre-trained LLMs to 256k with LongRoPE}. Using LLaMA2 as a case study, we perform a search aimed at achieving context window sizes of 128k and 256k. At this phase, the enlargement ratios correspond to 32$\times$ and 64$\times$, respectively.

$\Diamond$ \textit{Fine-tuning to 256k}. Next, we adapt the pre-trained LLM to handle a 256k context window through a staged fine-tuning process. Initially, LLaMA2 undergoes 400 fine-tuning steps utilizing RoPE rescaling factors corresponding to 128k. Once this phase is complete, we update the model by substituting the RoPE rescaling factors with those set for 256k on the obtained checkpoint, followed by an extra 600 steps of fine-tuning. This incremental approach demonstrates superior efficiency compared to immediately fine-tuning for the 256k context window size.

$\Diamond$ \textit{LongRoPE search enables further extending fine-tuned LLM context windows up to 2048k}. In the last stage, we apply an additional search procedure to the already fine-tuned LLM with a context window of 256k. Through this approach, we achieve an exceptional context window size of 2048k without any additional fine-tuning required. The overall extension corresponds to a $512\times$ increase.

\textbf{Shorter context window recovery}. Upon increasing the context window to a very large 2048k length, we observe that model performance deteriorates inside the standard context window. This phenomenon has been previously documented in relation to positional interpolation~\cite{pi}, which confines the position embeddings corresponding to the original context window to a significantly compressed subspace in higher-dimensional space. This compression adversely impacts language modeling accuracy. When extending by a factor of 512$\times$, the positions associated with the initial 4k context window are especially densely concentrated.

To address this issue, we conduct an additional evolution search on the extended LLM to fine-tune the RoPE rescale factors specifically for shorter context lengths (such as 4k and 8k). Since these lengths demand less positional interpolation, we constrain the upper bound of the searched $\lambda$ accordingly. At inference time, the LLM adaptively modifies the associated RoPE rescale factors.

 \section{Experiments}

\label{sec:exp}
\subsection{Setup}
\textbf{Evaluation Tasks and Models}. LongRoPE is integrated with both LLaMA2-7B and Mistral-7B models, whose capabilities are assessed along three dimensions: (1) the perplexity exhibited by extended-context LLMs when processing lengthy textual inputs; (2) performance on the Passkey retrieval task, which evaluates the capacity of a model to extract a basic passkey from a large volume of non-informative text; and (3) accuracy on established LLM benchmarks using a restricted context window of 4096 tokens.

\textbf{Fine-tuning}. For LLaMA2, a learning rate of $2 \times 10^{-5}$ with a linear decay schedule and a global batch size of 32 is applied. The model is fine-tuned over 400 steps using the Redpajama~\cite{redpajama} dataset, split into segments of 128k tokens and enclosed with BOS and EOS markers. Subsequently, starting from the obtained checkpoint, an extra 600 steps are conducted to reach a 256k context window. The training at the 128k context length employs 8 A100 GPUs via the distributed training system~\cite{cube}, whereas the 256k context window configuration utilizes 16 A100 GPUs. For Mistral, training utilizes a constant learning rate of $1 \times 10^{-6}$ alongside a global batch size of 64. Following the setup in YaRN~\cite{yarn}, both the 128k and 256k context models are fine-tuned for 400 steps on Together Computer's Long-Data Collections~\cite{mistraldata}, each with a sequence length of 16k, utilizing 4 A100 GPUs throughout.

\textbf{Search}. For target window sizes up to 256k, we set $P$=64, $N_1$=$N_2$=16, $p$=0.3, $\mathcal{T}$=40, and choose the top 32 candidates for mutation and crossover in each iteration. Perplexity is determined using 5 randomly selected samples from the PG19 validation dataset, ensuring that each meets the required minimum target context length. When dealing with windows larger than 512k, we reduce the population as well as the mutation and crossover sizes by half. In this case, perplexity is evaluated on 3 randomly chosen samples from the Pile-Books3~\cite{pile} validation dataset.

\textbf{Baselines}. To achieve a 2048k context window, we conducted fine-tuning on models with 128k and 256k context lengths. These procedures result in LongRoPE-2048k (ft=128k) and LongRoPE-2048k (ft=256k) for LLaMA2 and Mistral, respectively. We benchmark these four models against leading approaches for extending context windows, focusing on open-source LLMs that have been fine-tuned using positional interpolation techniques such as PI, NTK, and YaRN. The set of comparisons includes models like Together-32k~\cite{together}, Code LLaMA~\cite{codellama}, LongLoRA-full-FT-100k~\cite{longlora}, as well as YaRN-LLaMA and YaRN-Mistral~\cite{yarn}.

\subsection{Main Results}

\label{sec:mainresult}
\textbf{Long sequence language modeling within 256k}. Our initial experiments focus on benchmarking against leading extended LLMs with evaluation conducted up to a context length of 256k tokens. To illustrate the broad applicability of our approach, we employ two benchmark datasets: Proof-pile~\cite{pg19} and PG19~\cite{pile}, each evaluated on their respective test splits. Perplexity is measured at varying context lengths using a sliding window of size 256. For PG19, all 100 documents in the test set are included in the evaluation. For Proof-pile, aligning with the methodology of YaRN~\cite{yarn}, we randomly choose 10 samples, ensuring each selected sample contains at least 128k tokens.

Table~\ref{tbl:proof-pile} and Table~\ref{tbl:pg19} present a comparison of the perplexity values for LLaMA2 and Mistral models, both augmented with various interpolation strategies, on the Proof-pile and PG19 datasets, respectively. We emphasize two main findings: \textbf{(1)} our enhanced models consistently exhibit reduced perplexity as evaluation length increases from 4k up to 256k, indicating effective utilization of extended context. \textbf{(2)} Despite operating with a 16$\times$ longer context window—a scenario that typically compromises performance on shorter inputs—our LongRoPE-2048k models surpass leading baseline results within a 256k context window.

\textbf{Long sequence language modeling beyond 2000k}. To assess performance on exceptionally lengthy documents, we utilize the Books3~\cite{pile} dataset. For efficient evaluation, we randomly sample 20 books, each longer than 2048k tokens, and apply a sliding window approach with a size of 256k.

As detailed in Table~\ref{tbl:books3}, LongRoPE effectively broadens the context window of both LLaMA2-7B and Mistral-7B models to 2048k, while maintaining or surpassing baseline perplexity at shorter sequence lengths ranging from 8k to 128k. Clear distinctions arise between LLaMA2 and Mistral models at the 2048k length. Mistral demonstrates superior results compared to baselines at limited lengths; however, its perplexity exceeds 7 past the 256k mark. For LLaMA2, the observed results are consistent with expectations: its perplexity steadily decreases as the context grows longer, though there are slight upticks at the 1024k and 2048k marks. Additionally, for LLaMA2, LongRoPE-2048k yields better performance when fine-tuned at a window of 256k instead of 128k, attributed to the smaller ratio in the secondary extension (that is, 8$\times$ compared to 16$\times$). Conversely, for Mistral, optimal outcomes are achieved at a 128k fine-tuning length. This trend stems primarily from the decision to adopt a 16k training length for both Mistral's 128k and 256k fine-tuning—an approach based on YaRN’s configuration—which in turn impacts Mistral’s ability to expand its context window post fine-tuning.

\textbf{Passkey retrieval}. In this section, we investigate how the size of the effective context window impacts performance on generation tasks. To do this, we employ the synthetic passkey retrieval task introduced by~\cite{passkey}. In this setup, the model must locate a randomly chosen passkey, defined as a five-digit number, that is embedded somewhere within a lengthy document. The precise prompt format used is described in the appendix. For our experiments, we conduct 10 runs of the passkey retrieval task, each time inserting the passkey at a randomly selected position sampled uniformly throughout the evaluation context window.

Fig.~\ref{fig:passkey} presents a comparison of retrieval accuracy against several baselines. The accuracy of prior models drops precipitously to zero once the context size surpasses 128k tokens. Conversely, our LongRoPE-LLaMA2-2048k (ft=256k) demonstrates robust retrieval performance, achieving at least 90\% accuracy across the entire range from 4k to 2048k tokens—a notably challenging retrieval scenario at the million-token scale. LongRoPE-Mistral-2048k (ft=128k) consistently achieves perfect (100\%) accuracy up to 1800k tokens, before decreasing to 60\% at 2048k. This behavior is consistent with the findings in Table~\ref{tbl:books3}, where a modest increase in perplexity is also observed at 2048k.

\textbf{Evaluation on established benchmarks with the standard context length}. For the assessment of LongRoPE-2048k models within their original context window, we utilize the Hugging Face Open LLM Leaderboard~\cite{huggingfaceboard} in both zero-shot and few-shot scenarios. Specifically, the evaluation comprises 25-shot ARC-Challenge~\cite{allenai:arc}, 10-shot HellaSwag~\cite{zellers2019hellaswag}, 5-shot MMLU~\cite{mmlu}, and 0-shot TruthfulQA~\cite{lin2021truthfulqa}.

According to Table~\ref{tbl:commonsense}, our models attain performance similar to that of the original benchmark, which was created for a shorter context window. Notably, they surpass the original Mistral's results on TruthfulQA by an increase of +0.5\%. While LongRoPE-LLaMA2-2048k, after fine-tuning at 256k, exhibits a somewhat greater decline in performance, its results are still generally consistent with acceptable limits across the majority of evaluated tasks.

\subsection{Ablation Results}

\textbf{Effectiveness of the second positional interpolation}. Within the framework of our progressive extension approach, a second application of non-uniform positional interpolation is performed on the fine-tuned extended LLMs using our search-based method. To assess this approach, we conduct evaluations on our fine-tuned LLaMA2-256k model, further scaling it up to 512k, 1024k, and 2048k positions through both PI and YaRN methods. According to the results presented in Table~\ref{tbl:secondsearch}, our implementation of non-uniform positional interpolation maintains perplexity at a stable level. By contrast, as the extension ratio increases, PI and YaRN demonstrate a sharp rise in perplexity.

\textbf{Recovery performance at shorter context windows.} In order to address the decline in performance observed with shorter context lengths, we re-optimized the RoPE parameters for LongRoPE-2048k using our search method. In this process, we reduced the upper limit on scale factors, thereby reducing interpolation for shorter contexts of 4k and 8k tokens. Table~\ref{tbl:shortperformance} presents a perplexity comparison of LongRoPE-LLaMA2-2048k on Proof-pile at both 4k and 8k tokens, as well as corresponding average LLM benchmark accuracy. The findings highlight notable gains in performance when operating on shorter context lengths.

\textbf{Analysis on the two forms of non-uniformities}. To isolate the contributions of each non-uniformity, we conduct an ablation study. Specifically, we perform two sets of experiments: (i) for LLaMA2-7B, we extend sequence lengths to 16k and 32k using various approaches—PI, searching over the RoPE dimension exclusively, and searching over both types of non-uniformities; (ii) for our LLaMA2 model fine-tuned at a 256k sequence length, we extend to 2048k employing the same strategies. Perplexity scores are obtained without any further fine-tuning. As indicated in Table~\ref{tbl:searchdimension}, introducing non-uniformity in the RoPE dimension yields substantial perplexity reductions compared to the linear interpolation approach of PI. Incorporating non-uniformity in token position leads to clear improvements at 16k and 32k sequence lengths, whereas its influence diminishes at 2048k, likely due to the challenges posed by extremely long contexts. Retaining only the initial set of tokens without applying interpolation appears ineffective in such scenarios, which we designate as an area for future investigation.

\section{Related Works}

In addition to methods based on position interpolation, this section discusses related works of other approaches. 

\textbf{Retrieval-based} methods leverage an external memory to store extensive historical context and employ retrieval components to obtain relevant documents during inference~\cite{longllama,wang2023augmenting,borgeaud2022improving}. Such strategies generally require direct architectural adjustments to LLMs. By contrast, our approach adopts a more streamlined strategy, introducing only minimal changes to the positional embeddings. Additionally, we demonstrate the flexibility of our model to address a wider range of long-context applications beyond simple retrieval, including long document summarization and few-shot learning.

\textbf{Attention-based context window extensions}. In addition to positional embedding interpolation, other studies have focused on lengthening the input context within the constraints of the standard LLM context window by altering attention mechanisms~\cite{han2023lminfinite, streamingllm,parallelwindow}. Central to these approaches is the use of innovative attention masks to alleviate the problem of attention explosion that arises with the inclusion of new token positions. These techniques can be effectively combined with positional interpolation strategies, as they address different aspects of context extension.

\textbf{Fine-tuning based approaches} address the challenge of adapting pre-trained LLMs to longer sequences by modifying their position embeddings. Approaches such as Code LLaMA~\cite{codellama}, LLaMA2 Long~\cite{xiong2023effective}, and ScaledRoPE~\cite{liu2023scaling} implement a substantially increased RoPE base value, followed by fine-tuning at the desired sequence length. The method we propose is applicable across a range of target lengths and enables processing of lengths exceeding 2M tokens. As expanding fine-tuning to longer contexts (greater than 128k tokens) is computationally intensive and requires significant GPU capabilities, recent techniques like LongLoRA~\cite{longlora} and PoSE~\cite{zhu2023pose} have been introduced to address this scalability concern. Our approach can be applied in conjunction with these efficient fine-tuning strategies.

\section{Conclusion}

In this paper, we introduce LongRoPE, a technique that significantly increases the context window of LLMs up to an unprecedented 2048k tokens, all while preserving their performance within the original, more limited context size. Our approach leverages two distinct types of non-uniformities inherent in RoPE positional encodings, employing an efficient evolutionary search to do so. This method yields two primary advantages: it facilitates effective initialization for fine-tuning, and it supports an 8$\times$ expansion of the context window without requiring fine-tuning. Building on this foundation, we develop a progressive extension framework, wherein 256k-token fine-tuned LLMs are used to further extend the context window up to 2048k tokens, again without needing additional fine-tuning steps. Comprehensive experiments confirm the robustness and effectiveness of LongRoPE. We anticipate that our LongRoPE-2048k models will unlock a variety of novel long-context use cases and drive future developments in the field.

\section*{Broader Impacts} The objective of this paper is to contribute to progress within the discipline of Machine Learning. While our research may carry various implications for society, we do not identify any particular issues that require special attention in this context.

	\balance

\end{document}